%-------------------------
% Cover letter - Veeva Systems, Senior Software Engineer, Backend (Veeva Ostro)
% Author : Aymen Daoudi
% Reuses Parameters.sty / Commands.sty from the resume so both documents match.
%------------------------

\documentclass[letterpaper,10pt]{article}

% Parameters.sty, Commands.sty and Images/ live in the repo root, one level up.
% \input@path lets \usepackage resolve them from inside this folder.
\makeatletter
\def\input@path{{../}}
\makeatother

\usepackage{Parameters}
\usepackage{Commands}

% graphicx is loaded by Parameters.sty, so \graphicspath must come after it.
% This resolves the header icons under Images/.
\graphicspath{{../}}

% Improves inter-word spacing in justified text; keeps the rivers down.
\usepackage{microtype}

% Letter geometry: the resume's 1cm margins are too tight for running prose.
\geometry{left=2.6cm, right=2.6cm, top=1.6cm, bottom=1.6cm}

% Block paragraph style, no indentation.
\setlength{\parindent}{0pt}
\setlength{\parskip}{7pt}
\setlength{\footskip}{12pt}
\linespread{1.05}

\begin{document}

% \roboto is a declaration, not a wrapper. Commands.sty issues it per-macro, so
% the resume renders in Roboto; running prose here needs it set once, document
% wide, or it falls back to the Latin Modern serif default.
\roboto

% Parameters.sty issues a global \raggedright for the resume layout. Undo it so
% the letter body is justified. \parindent stays at 0pt for block paragraphs.
\rightskip=0pt
\leftskip=0pt
\parfillskip=0pt plus 1fil

%----------HEADING (mirrors Parts/Heading.tex)-----------------

\begin{tabular*}{\textwidth}{l@{\extracolsep{\fill}}l}
  \multirow{2}{*}{\Name{\huge A\LARGE YMEN \huge D\LARGE AOUDI}}
    & \Addresse{\quad Montréal, Québec, Canada} \\
    & \EmailAddresse{\quad \href{mailto:iaymendaoudi@gmail.com}{iaymendaoudi@gmail.com}} \\
  \RoleDescription{Software Development Engineer}
    & \PhoneNumbers{\quad +1 438 927 9580}{} \\
    & \faLinkedin \enspace \enspace {\url{https://www.linkedin.com/in/aymendaoudi}}
\end{tabular*}

\vspace{6pt}
{\color{black}\rule{\textwidth}{0.4pt}}
\vspace{10pt}

%----------LETTER-----------------

{\footnotesize August 22, 2026}

\vspace{2pt}

Dear Veeva Hiring Team,

I am applying for the Senior Software Engineer, Backend position with Veeva
Ostro. It reads as a close match to what I do now and want to keep doing: Python
and Django at depth, backend services carrying real load, and enough weight given
to clean architecture and testing that code health counts as part of delivery
rather than cleanup afterwards. That the platform serves life sciences makes it
more appealing rather than less, since the consequences of the work are easy to
see.

I am a Senior Software Engineer at Numa, working in Python across Django, FastAPI
and Flask on the services behind an AI phone assistant for car dealerships, over
Postgres. Most of my time goes to the backend integration layer between our
platform and the dealer management systems it depends on, including the Reynolds
\& Reynolds integration I led. The rest goes to keeping it healthy in production:
watching live traffic, tracing failures back through the stack, and fixing them
at the source rather than the symptom. Python and Django have been part of my
work since 2020, with Postgres alongside them.

Before Numa I spent three years at Divisions Maintenance Group building the
backend behind a platform connecting maintenance providers with work across the
United States. Earlier roles took me through maritime compliance software at
Lloyd's Register OneOcean, electric vehicle charging at AddEnergie, and a
financial policy and insurance platform at Stratec Global, and my first years
were spent building internal reporting systems for a pharmaceutical group. The
work has been much the same throughout: designing backend systems that have to
keep running, and pushing quality upstream through code review, testing and
written design. My MSc in Computer Science is from UQAM, with research published
at SAC'19.

What draws me to Veeva Ostro is that the platform sits between the companies
making treatments and the people who need them, which puts a premium on exactly
what the role asks for: uptime, clean architecture and low technical debt. Those
are easy to claim and harder to keep to when delivery is quick, and that is where
I would contribute, taking features from design through production, giving
colleagues the review I would want on my own work, and leaving the codebase
easier to work in than I found it. I am based in Montréal and have worked
remotely with distributed teams since 2022, so Work Anywhere suits how I already
operate.

Thank you for considering my application. I would be glad to discuss how I can
contribute to the Ostro team.

\vspace{6pt}

Respectfully,

\vspace{4pt}

Aymen Daoudi

\end{document}
